\documentclass[11pt, a4paper]{article}
\usepackage[a4paper, top=2.5cm, bottom=2.5cm, left=2cm, right=2cm]{geometry}
\usepackage{fontspec}
\usepackage[english, bidi=basic, provide=*]{babel}
\babelprovide[import, onchar=ids fonts]{english}
\babelfont{rm}{TeX Gyre Termes}
\babelfont{sf}{TeX Gyre Heros}
\babelfont{tt}{TeX Gyre Cursor}
\usepackage{enumitem}
\usepackage{amsmath}

\title{MCQ Practice Set: The Cell Cycle}
\author{Subject: Biology (Chapter 12)}
\date{}

\begin{document}
\maketitle

\section*{Practice Questions}

\begin{enumerate}[leftmargin=*, label=\arabic*.]

\item \textbf{A cell's endowment of DNA, its genetic information, is called its:}
\begin{enumerate}[label=\Alph*.]
\item Chromatin
\item Genome
\item Chromatid
\item Centromere
\item Nucleosome
\end{enumerate}
\textbf{Answer:} B \\
\textbf{Explanation:} The genome represents the complete set of genetic material in a cell.

\item \textbf{How many chromosomes are in a typical human somatic cell?}
\begin{enumerate}[label=\Alph*.]
\item 23
\item 44
\item 46
\item 48
\item 92
\end{enumerate}
\textbf{Answer:} C \\
\textbf{Explanation:} Human somatic cells (body cells) contain 46 chromosomes, consisting of two sets of 23.

\item \textbf{Human gametes, such as sperm or eggs, have \underline{\hspace{1cm}} chromosomes.}
\begin{enumerate}[label=\Alph*.]
\item 23
\item 46
\item 92
\item 12
\item 48
\end{enumerate}
\textbf{Answer:} A \\
\textbf{Explanation:} Gametes are haploid cells and contain half the number of chromosomes found in somatic cells.

\item \textbf{The complex of DNA and proteins that is the building material of eukaryotic chromosomes is called:}
\begin{enumerate}[label=\Alph*.]
\item Cohesin
\item Centromere
\item Chromatin
\item Genome
\item Spindle
\end{enumerate}
\textbf{Answer:} C \\
\textbf{Explanation:} Chromatin is the mass of DNA and protein that condenses to form chromosomes during cell division.

\item \textbf{Identical joined copies of a single duplicated chromosome are called:}
\begin{enumerate}[label=\Alph*.]
\item Daughter chromosomes
\item Homologous pairs
\item Sister chromatids
\item Centrioles
\item Asters
\end{enumerate}
\textbf{Answer:} C \\
\textbf{Explanation:} Sister chromatids are the two identical copies held together by cohesins after DNA replication.

\item \textbf{The region of a chromosome where two sister chromatids are most closely attached is the:}
\begin{enumerate}[label=\Alph*.]
\item Kinetochore
\item Centrosome
\item Centromere
\item Cleavage furrow
\item Metaphase plate
\end{enumerate}
\textbf{Answer:} C \\
\textbf{Explanation:} The centromere is the narrow "waist" of the duplicated chromosome where chromatids are linked.

\item \textbf{The division of the genetic material in the nucleus is called:}
\begin{enumerate}[label=\Alph*.]
\item Cytokinesis
\item Meiosis
\item Mitosis
\item Binary fission
\item Interphase
\end{enumerate}
\textbf{Answer:} C \\
\textbf{Explanation:} Mitosis specifically refers to the nuclear division that results in identical daughter nuclei.

\item \textbf{The division of the cytoplasm, which usually follows mitosis, is called:}
\begin{enumerate}[label=\Alph*.]
\item Telophase
\item Anaphase
\item Cytokinesis
\item S phase
\item G2 phase
\end{enumerate}
\textbf{Answer:} C \\
\textbf{Explanation:} Cytokinesis is the physical separation of the cytoplasm into two distinct daughter cells.

\item \textbf{Which of the following is NOT part of interphase?}
\begin{enumerate}[label=\Alph*.]
\item G1 phase
\item S phase
\item G2 phase
\item Mitotic phase
\item Metabolic activity
\end{enumerate}
\textbf{Answer:} D \\
\textbf{Explanation:} Interphase includes G1, S, and G2; the Mitotic (M) phase is a separate part of the cell cycle.

\item \textbf{DNA replication (synthesis) occurs entirely during which phase?}
\begin{enumerate}[label=\Alph*.]
\item G1 phase
\item S phase
\item G2 phase
\item Prophase
\item Metaphase
\end{enumerate}
\textbf{Answer:} B \\
\textbf{Explanation:} The S phase stands for "synthesis," where the cell copies its chromosomes.

\item \textbf{Which phase of the cell cycle is typically the shortest?}
\begin{enumerate}[label=\Alph*.]
\item G1 phase
\item S phase
\item G2 phase
\item M phase
\item Interphase
\end{enumerate}
\textbf{Answer:} D \\
\textbf{Explanation:} The mitotic (M) phase, including mitosis and cytokinesis, is usually the shortest part of the cycle.

\item \textbf{In a 24-hour human cell cycle, the M phase would likely occupy approximately how much time?}
\begin{enumerate}[label=\Alph*.]
\item 10-12 hours
\item 5-6 hours
\item Less than 1 hour
\item 20 hours
\item 12 hours
\end{enumerate}
\textbf{Answer:} C \\
\textbf{Explanation:} According to the text, M phase is the shortest, often taking less than one hour in a 24-hour cycle.

\item \textbf{The five stages of mitosis, in order, are:}
\begin{enumerate}[label=\Alph*.]
\item Prophase, Metaphase, Anaphase, Telophase, Cytokinesis
\item Prophase, Prometaphase, Metaphase, Anaphase, Telophase
\item G1, S, G2, M, C
\item Interphase, Prophase, Metaphase, Anaphase, Telophase
\item Prophase, Prometaphase, Metaphase, Telophase, Anaphase
\end{enumerate}
\textbf{Answer:} B \\
\textbf{Explanation:} The five conventional stages are prophase, prometaphase, metaphase, anaphase, and telophase.

\item \textbf{During which stage do the nucleoli disappear and the chromatin fibers condense?}
\begin{enumerate}[label=\Alph*.]
\item Interphase
\item Prophase
\item Metaphase
\item Anaphase
\item Telophase
\end{enumerate}
\textbf{Answer:} B \\
\textbf{Explanation:} Prophase is characterized by the visible condensation of chromosomes and the disappearance of nucleoli.

\item \textbf{What is the "mitotic spindle" composed of?}
\begin{enumerate}[label=\Alph*.]
\item DNA and histones
\item Actin and myosin
\item Microtubules and associated proteins
\item Phospholipids and proteins
\item Keratin fibers
\end{enumerate}
\textbf{Answer:} C \\
\textbf{Explanation:} The spindle is a microtubule-based structure that moves chromosomes during mitosis.

\item \textbf{The assembly of spindle microtubules starts at the \underline{\hspace{1cm}} in animal cells.}
\begin{enumerate}[label=\Alph*.]
\item Centromere
\item Centrosome
\item Kinetochore
\item Nucleolus
\item Cleavage furrow
\end{enumerate}
\textbf{Answer:} B \\
\textbf{Explanation:} The centrosome acts as the microtubule-organizing center (MTOC) for the mitotic spindle.

\item \textbf{An "aster" is a:}
\begin{enumerate}[label=\Alph*.]
\item Structure that helps pinch the cell in two
\item Radial array of short microtubules extending from the centrosome
\item Specialized protein structure at the centromere
\item Plane equidistant between the spindle's two poles
\item Type of enzyme that cleaves cohesins
\end{enumerate}
\textbf{Answer:} B \\
\textbf{Explanation:} Asters are the star-shaped radial arrays of microtubules around the centrosomes.

\item \textbf{A specialized protein structure at the centromere that attaches to spindle microtubules is the:}
\begin{enumerate}[label=\Alph*.]
\item Aster
\item Centriole
\item Kinetochore
\item Metaphase plate
\item Separase
\end{enumerate}
\textbf{Answer:} C \\
\textbf{Explanation:} Each sister chromatid has a kinetochore where spindle fibers attach.

\item \textbf{At metaphase, chromosomes align on an imaginary plane called the:}
\begin{enumerate}[label=\Alph*.]
\item Cleavage furrow
\item Metaphase plate
\item Cell plate
\item Origin of replication
\item Nuclear envelope
\end{enumerate}
\textbf{Answer:} B \\
\textbf{Explanation:} The metaphase plate is the midway plane between the spindle's two poles.

\item \textbf{Which enzyme is responsible for cleaving the cohesins that hold sister chromatids together?}
\begin{enumerate}[label=\Alph*.]
\item Polymerase
\item Separase
\item Kinase
\item Myosin
\item Actin
\end{enumerate}
\textbf{Answer:} B \\
\textbf{Explanation:} Separase initiates anaphase by breaking the protein bonds between sister chromatids.

\item \textbf{Anaphase is characterized by:}
\begin{enumerate}[label=\Alph*.]
\item The reformation of the nuclear envelope
\item Chromosomes aligning at the center
\item The sudden parting of sister chromatids
\item DNA replication
\item The disappearance of the nucleolus
\end{enumerate}
\textbf{Answer:} C \\
\textbf{Explanation:} Anaphase begins when cohesins are cleaved and sister chromatids move toward opposite poles.

\item \textbf{According to the "Pac-man" mechanism, chromosomes move toward the poles because:}
\begin{enumerate}[label=\Alph*.]
\item They are reeled in by motor proteins at the poles
\item Motor proteins on kinetochores "walk" along microtubules that depolymerize at the kinetochore end
\item The cell's cytoplasm flows toward the poles
\item Actin and myosin contract
\item The spindle fibers simply push them
\end{enumerate}
\textbf{Answer:} B \\
\textbf{Explanation:} This mechanism involves the depolymerization of microtubules at the kinetochore as the motor proteins pass.

\item \textbf{In animal cells, cytokinesis involves the formation of a:}
\begin{enumerate}[label=\Alph*.]
\item Cell plate
\item Cell wall
\item Cleavage furrow
\item Spindle pole
\item Nucleolus
\end{enumerate}
\textbf{Answer:} C \\
\textbf{Explanation:} A cleavage furrow forms in animal cells, which is a shallow groove pinched by a contractile ring.

\item \textbf{The contractile ring that forms the cleavage furrow is made of:}
\begin{enumerate}[label=\Alph*.]
\item Tubulin and motor proteins
\item Actin microfilaments and myosin
\item DNA and histones
\item Cellulose and pectins
\item Keratin and intermediate filaments
\end{enumerate}
\textbf{Answer:} B \\
\textbf{Explanation:} Actin and myosin interact to contract the ring and pinch the cell in two.

\item \textbf{In plant cells, cytokinesis occurs through the formation of a:}
\begin{enumerate}[label=\Alph*.]
\item Cleavage furrow
\item Cell plate
\item Centrosome
\item Origin of replication
\item Binary fission ring
\end{enumerate}
\textbf{Answer:} B \\
\textbf{Explanation:} Plant cells build a cell plate from Golgi-derived vesicles because their rigid cell walls prevent furrowing.

\item \textbf{Binary fission is the method of cell division used by:}
\begin{enumerate}[label=\Alph*.]
\item All eukaryotes
\item Plants and animals
\item Bacteria and archaea
\item Fungi only
\item Viruses
\end{enumerate}
\textbf{Answer:} C \\
\textbf{Explanation:} Binary fission is asexual reproduction in prokaryotes (bacteria and archaea).

\item \textbf{The specific place on a bacterial chromosome where replication begins is called the:}
\begin{enumerate}[label=\Alph*.]
\item Centromere
\item Kinetochore
\item Origin of replication
\item Metaphase plate
\item Telomere
\end{enumerate}
\textbf{Answer:} C \\
\textbf{Explanation:} Replication starts at the origin of replication and proceeds in both directions.

\item \textbf{What is a key difference between binary fission in bacteria and mitosis in eukaryotes?}
\begin{enumerate}[label=\Alph*.]
\item Bacteria do not replicate their DNA
\item Mitosis involves a mitotic spindle; binary fission does not
\item Binary fission produces four daughter cells
\item Mitosis only occurs in single-celled organisms
\item Bacteria do not use proteins during division
\end{enumerate}
\textbf{Answer:} B \\
\textbf{Explanation:} Bacteria lack a visible mitotic spindle and microtubules for chromosome movement.

\item \textbf{Which group of organisms shows intermediate stages of cell division evolution where the nuclear envelope remains intact?}
\begin{enumerate}[label=\Alph*.]
\item Bacteria and archaea
\item Most plants and animals
\item Dinoflagellates and diatoms
\item Viruses
\item Mammals
\end{enumerate}
\textbf{Answer:} C \\
\textbf{Explanation:} In dinoflagellates and diatoms, the nucleus divides without the envelope breaking down.

\item \textbf{What is a "checkpoint" in the cell cycle?}
\begin{enumerate}[label=\Alph*.]
\item A physical barrier in the cytoplasm
\item A control point where stop and go-ahead signals regulate the cycle
\item The moment DNA replication finishes
\item The start of prophase
\item The time when a cell dies
\end{enumerate}
\textbf{Answer:} B \\
\textbf{Explanation:} Checkpoints allow the cell to verify that crucial processes have been completed before proceeding.

\item \textbf{The three most important checkpoints are found in which phases?}
\begin{enumerate}[label=\Alph*.]
\item Prophase, Metaphase, Anaphase
\item G1, S, G2
\item G1, G2, M
\item S, G2, M
\item Interphase, Mitosis, Cytokinesis
\end{enumerate}
\textbf{Answer:} C \\
\textbf{Explanation:} The G1, G2, and M checkpoints are the primary regulatory gates.

\item \textbf{If a cell does not receive a go-ahead signal at the G1 checkpoint, it may enter a non-dividing state called:}
\begin{enumerate}[label=\Alph*.]
\item G2 phase
\item S phase
\item G0 phase
\item M phase
\item Cytokinesis
\item 
\end{enumerate}
\textbf{Answer:} C \\
\textbf{Explanation:} Most cells in the human body are in the G0 phase, a resting or non-dividing state.

\item \textbf{The cell cycle control system is driven by regulatory proteins called:}
\begin{enumerate}[label=\Alph*.]
\item Actins and myosins
\item Cyclins and Cdks
\item Histones and cohesins
\item Tubulins and dyneins
\item Separases and polymerases
\end{enumerate}
\textbf{Answer:} B \\
\textbf{Explanation:} Cyclins and Cyclin-dependent kinases (Cdks) are the primary pacing molecules of the cycle.

\item \textbf{What does "MPF" stand for in the context of cell cycle regulation?}
\begin{enumerate}[label=\Alph*.]
\item Mitosis Protein Factor
\item Maturation-promoting factor
\item Metabolic Process Fiber
\item Microtubule Polymerization Fragment
\item Major Protein Fusion
\end{enumerate}
\textbf{Answer:} B \\
\textbf{Explanation:} MPF triggers the cell's passage into the M phase past the G2 checkpoint.

\item \textbf{The activity of a Cdk rises and falls based on the concentration of its \underline{\hspace{1cm}} partner.}
\begin{enumerate}[label=\Alph*.]
\item ATP
\item Glucose
\item Cyclin
\item Actin
\item DNA
\end{enumerate}
\textbf{Answer:} C \\
\textbf{Explanation:} Cdks require cyclin binding to become active; cyclin levels fluctuate rhythmically.

\item \textbf{A kinase is an enzyme that:}
\begin{enumerate}[label=\Alph*.]
\item Breaks down proteins
\item Replicates DNA
\item Activates or inactivates other proteins by phosphorylating them
\item Synthesizes lipids
\item Transports oxygen
\end{enumerate}
\textbf{Answer:} C \\
\textbf{Explanation:} Kinases use phosphate groups to modify the activity of target proteins.

\item \textbf{Which stage of the cell cycle involves metabolic activity, growth, and preparation for division?}
\begin{enumerate}[label=\Alph*.]
\item S phase
\item M phase
\item G2 phase
\item Anaphase
\item Cytokinesis
\end{enumerate}
\textbf{Answer:} C \\
\textbf{Explanation:} G2 is the "second gap" phase where final preparations for M phase occur.

\item \textbf{Walther Flemming coined which terms in 1882?}
\begin{enumerate}[label=\Alph*.]
\item DNA and RNA
\item Mitosis and chromatin
\item Genome and gene
\item Centromere and centrosome
\item Prokaryote and eukaryote
\end{enumerate}
\textbf{Answer:} B \\
\textbf{Explanation:} Flemming was an anatomist who first observed the behavior of chromosomes.

\item \textbf{A typical human cell has about \underline{\hspace{1cm}} of DNA.}
\begin{enumerate}[label=\Alph*.]
\item 2 mm
\item 2 cm
\item 2 m
\item 2 km
\item 2 $\mu$m
\end{enumerate}
\textbf{Answer:} C \\
\textbf{Explanation:} The text states that a typical human cell has about 2 meters of DNA.

\item \textbf{Meiosis produces daughter cells with how many sets of chromosomes?}
\begin{enumerate}[label=\Alph*.]
\item Two sets
\item Four sets
\item One set
\item No sets
\item Three sets
\end{enumerate}
\textbf{Answer:} C \\
\textbf{Explanation:} Meiosis reduces the chromosome number by half to produce haploid gametes.

\item \textbf{The proteins that hold sister chromatids together all along their lengths are called:}
\begin{enumerate}[label=\Alph*.]
\item Histones
\item Tubulins
\item Cohesins
\item Actins
\item Myosins
\end{enumerate}
\textbf{Answer:} C \\
\textbf{Explanation:} Cohesins mediate sister chromatid cohesion.

\item \textbf{What happens to the sister chromatids after they separate in anaphase?}
\begin{enumerate}[label=\Alph*.]
\item They are called daughter chromosomes
\item They dissolve
\item They leave the cell
\item They turn back into chromatin fibers immediately
\item They fuse with the plasma membrane
\end{enumerate}
\textbf{Answer:} A \\
\textbf{Explanation:} Once separated, each chromatid is considered an individual chromosome.

\item \textbf{Which phase of mitosis involves the nuclear envelope fragments and microtubules invading the nuclear area?}
\begin{enumerate}[label=\Alph*.]
\item Prophase
\item Prometaphase
\item Metaphase
\item Anaphase
\item Telophase
\end{enumerate}
\textbf{Answer:} B \\
\textbf{Explanation:} During prometaphase, the envelope breaks down, allowing spindle access to chromosomes.

\item \textbf{Nonkinetochore microtubules are responsible for:}
\begin{enumerate}[label=\Alph*.]
\item Moving chromosomes
\item Replicating DNA
\item Elongating the whole cell during anaphase
\item Forming the nucleolus
\item Cleaving cohesins
\end{enumerate}
\textbf{Answer:} C \\
\textbf{Explanation:} They overlap and push against each other to stretch the cell.

\item \textbf{The daughter nuclei form during which stage?}
\begin{enumerate}[label=\Alph*.]
\item Prophase
\item Metaphase
\item Anaphase
\item Telophase
\item S phase
\end{enumerate}
\textbf{Answer:} D \\
\textbf{Explanation:} Telophase is the final nuclear stage where envelopes reform around the new sets of chromosomes.

\item \textbf{What is the "metaphase plate"?}
\begin{enumerate}[label=\Alph*.]
\item A solid structure in the center of the cell
\item An imaginary plane midway between the spindle poles
\item A layer of the cell wall
\item The site of DNA synthesis
\item The nucleolus
\end{enumerate}
\textbf{Answer:} B \\
\textbf{Explanation:} It is an imaginary plate used to describe chromosome alignment.

\item \textbf{In which phase do nucleoli reappear?}
\begin{enumerate}[label=\Alph*.]
\item Prophase
\item Metaphase
\item Anaphase
\item Telophase
\item G1 phase
\end{enumerate}
\textbf{Answer:} D \\
\textbf{Explanation:} Nucleoli reappear as the nuclei complete their formation in telophase.

\item \textbf{Which structure is NOT found in plant cells during mitosis?}
\begin{enumerate}[label=\Alph*.]
\item Chromosomes
\item Centrioles
\item Spindle
\item Nucleus
\item Cell plate
\end{enumerate}
\textbf{Answer:} B \\
\textbf{Explanation:} Centrioles are absent in plants, though they still form a spindle.

\item \textbf{E. coli chromosome is about \underline{\hspace{1cm}} times as long as the cell itself.}
\begin{enumerate}[label=\Alph*.]
\item 10
\item 100
\item 500
\item 1,000
\item 50
\end{enumerate}
\textbf{Answer:} C \\
\textbf{Explanation:} It is about 500 times longer, requiring extreme coiling and folding.

\item \textbf{The cell cycle control system was first studied using which organisms?}
\begin{enumerate}[label=\Alph*.]
\item Bacteria
\item Fruit flies
\item Mammalian cells in culture and yeasts
\item Plants
\item Viruses
\end{enumerate}
\textbf{Answer:} C \\
\textbf{Explanation:} Cell fusion experiments in mammals and genetic studies in yeast provided early evidence.

\item \textbf{During anaphase, chromosomes are pulled ahead of their \underline{\hspace{1cm}}.}
\begin{enumerate}[label=\Alph*.]
\item Centrosomes
\item Spindle poles
\item Chromatid arms
\item Nuclei
\item Origins
\end{enumerate}
\textbf{Answer:} C \\
\textbf{Explanation:} Because microtubules attach at the centromere, the centromeres lead the movement.

\item \textbf{The M checkpoint ensures that:}
\begin{enumerate}[label=\Alph*.]
\item DNA is replicated
\item The cell has grown enough
\item All kinetochores are properly attached to spindle microtubules
\item Cohesins are synthesized
\item ATP is available
\end{enumerate}
\textbf{Answer:} C \\
\textbf{Explanation:} This prevents anaphase from starting until all chromosomes are aligned and ready.

\item \textbf{How many chromatid arms does a duplicated chromosome have?}
\begin{enumerate}[label=\Alph*.]
\item One
\item Two
\item Three
\item Four
\item Eight
\end{enumerate}
\textbf{Answer:} D \\
\textbf{Explanation:} Each duplicated chromosome has two sister chromatids, and each chromatid has two arms (one on each side of the centromere).

\item \textbf{Which stage of mitosis is typically the shortest?}
\begin{enumerate}[label=\Alph*.]
\item Prophase
\item Metaphase
\item Anaphase
\item Telophase
\item Prometaphase
\end{enumerate}
\textbf{Answer:} C \\
\textbf{Explanation:} Anaphase often lasts only a few minutes.

\item \textbf{The rate of chromosome movement during anaphase is about:}
\begin{enumerate}[label=\Alph*.]
\item 1 mm/min
\item 1 $\mu$m/min
\item 1 cm/min
\item 1 m/min
\item 1 nm/min
\end{enumerate}
\textbf{Answer:} B \\
\textbf{Explanation:} Chromosomes move relatively slowly at about 1 micrometer per minute.

\item \textbf{What is the function of the "Pac-man" mechanism?}
\begin{enumerate}[label=\Alph*.]
\item To synthesize new DNA
\item To depolymerize microtubules and move chromosomes
\item To create the cell plate
\item To destroy the nuclear envelope
\item To replicate centrioles
\end{enumerate}
\textbf{Answer:} B \\
\textbf{Explanation:} It is one of the mechanisms for moving chromosomes toward the poles.

\item \textbf{Which proteins are related to eukaryotic tubulin in bacteria?}
\begin{enumerate}[label=\Alph*.]
\item FtsZ
\item MreB
\item Cohesin
\item Myosin
\item Actin
\end{enumerate}
\textbf{Answer:} A \\
\textbf{Explanation:} (Note: while not explicitly named "FtsZ" in the snippet, the text refers to "tubulin-like" proteins involved in pinching the membrane).

\item \textbf{Which phase follows G2?}
\begin{enumerate}[label=\Alph*.]
\item S phase
\item G1 phase
\item M phase
\item G0 phase
\item Prophase only
\end{enumerate}
\textbf{Answer:} C \\
\textbf{Explanation:} The mitotic phase follows the G2 preparation phase.

\item \textbf{What is the genome of a prokaryotic organism usually?}
\begin{enumerate}[label=\Alph*.]
\item Multiple linear DNA molecules
\item A single circular DNA molecule
\item RNA only
\item Proteins
\item A single linear strand
\end{enumerate}
\textbf{Answer:} B \\
\textbf{Explanation:} Most bacteria have a single circular chromosome.

\item \textbf{Condensation of chromosomes is necessary to:}
\begin{enumerate}[label=\Alph*.]
\item Make them long and thin
\item Allow them to be easily sorted and moved
\item Synthesize mRNA
\item Replicate DNA
\item Enter G0 phase
\end{enumerate}
\textbf{Answer:} B \\
\textbf{Explanation:} Large amounts of DNA must be packaged to be manageable during division.

\item \textbf{A hedgehog cell has 90 chromosomes. How many are in its gametes?}
\begin{enumerate}[label=\Alph*.]
\item 180
\item 45
\item 90
\item 23
\item 148
\end{enumerate}
\textbf{Answer:} B \\
\textbf{Explanation:} Gametes have half the number of somatic chromosomes.

\item \textbf{Which of these does NOT happen during interphase?}
\begin{enumerate}[label=\Alph*.]
\item Production of proteins
\item Organelle duplication
\item Spindle formation
\item DNA synthesis
\item Cell growth
\end{enumerate}
\textbf{Answer:} C \\
\textbf{Explanation:} The spindle begins to form during prophase of the mitotic phase.

\item \textbf{The G1 phase is also called the:}
\begin{enumerate}[label=\Alph*.]
\item Synthesis phase
\item Second gap
\item First gap
\item Mitotic phase
\item Resting phase
\end{enumerate}
\textbf{Answer:} C \\
\textbf{Explanation:} G1 stands for "first gap."

\item \textbf{Which part of interphase is most variable in length?}
\begin{enumerate}[label=\Alph*.]
\item G1 phase
\item S phase
\item G2 phase
\item M phase
\item Cytokinesis
\end{enumerate}
\textbf{Answer:} A \\
\textbf{Explanation:} G1 duration varies significantly depending on cell type.

\item \textbf{During prometaphase, "kinetochore microtubules" do what to chromosomes?}
\begin{enumerate}[label=\Alph*.]
\item Hold them perfectly still
\item Jerk them back and forth
\item Replicate them
\item Melt them
\item Hide them
\end{enumerate}
\textbf{Answer:} B \\
\textbf{Explanation:} The tug-of-war between opposite poles results in erratic movement until alignment.

\item \textbf{Which stage is the shortest in MITOSIS (not the whole cell cycle)?}
\begin{enumerate}[label=\Alph*.]
\item Prophase
\item Metaphase
\item Anaphase
\item Telophase
\item Prometaphase
\end{enumerate}
\textbf{Answer:} C \\
\textbf{Explanation:} Anaphase is the shortest stage of mitosis.

\item \textbf{The division of the cytoplasm is usually well under way by:}
\begin{enumerate}[label=\Alph*.]
\item Prophase
\item Metaphase
\item Anaphase
\item Late telophase
\item Early interphase
\end{enumerate}
\textbf{Answer:} D \\
\textbf{Explanation:} Cytokinesis overlaps with the end of mitosis.

\item \textbf{Polymerization refers to \underline{\hspace{1cm}}, while depolymerization refers to \underline{\hspace{1cm}}.}
\begin{enumerate}[label=\Alph*.]
\item Shortening; elongating
\item Melting; freezing
\item Elongating; shortening
\item Replicating; dividing
\item Condensing; expanding
\end{enumerate}
\textbf{Answer:} C \\
\textbf{Explanation:} Microtubules elongate by adding subunits (poly) and shorten by losing them (depoly).

\item \textbf{What provides the material to construct the mitotic spindle?}
\begin{enumerate}[label=\Alph*.]
\item The nucleus
\item Partially disassembled cytoskeleton
\item Mitochondria
\item The cell wall
\item Ribosomes
\end{enumerate}
\textbf{Answer:} B \\
\textbf{Explanation:} Cytoskeletal microtubules disassemble to provide tubulin for the spindle.

\item \textbf{Centrioles are \underline{\hspace{1cm}} for cell division to occur.}
\begin{enumerate}[label=\Alph*.]
\item Absolutely required in all cells
\item Not essential even in animal cells (spindle forms without them)
\item The only source of DNA
\item Found only in plants
\item Made of actin
\end{enumerate}
\textbf{Answer:} B \\
\textbf{Explanation:} Laser experiments showed spindles form even if centrioles are destroyed.

\item \textbf{What is the result of a draw in the "tug-of-war" during prometaphase?}
\begin{enumerate}[label=\Alph*.]
\item The cell dies
\item The chromosomes align at the metaphase plate
\item The chromosomes are ripped in half
\item The spindle collapses
\item DNA is destroyed
\end{enumerate}
\textbf{Answer:} B \\
\textbf{Explanation:} Equal tension from both poles centers the chromosomes.

\item \textbf{Separase becomes active when:}
\begin{enumerate}[label=\Alph*.]
\item The cell enters G1
\item All kinetochores are attached and aligned
\item DNA starts replicating
\item The cell plate forms
\item The nucleolus disappears
\end{enumerate}
\textbf{Answer:} B \\
\textbf{Explanation:} This is a critical regulatory step for starting anaphase.

\item \textbf{Motor proteins at the spindle poles \underline{\hspace{1cm}} chromosomes.}
\begin{enumerate}[label=\Alph*.]
\item Walk along
\item Reel in
\item Push
\item Duplicate
\item Ignore
\end{enumerate}
\textbf{Answer:} B \\
\textbf{Explanation:} This is an alternative/complementary mechanism to the Pac-man model.

\item \textbf{Nonkinetochore microtubules from opposite poles \underline{\hspace{1cm}} at the metaphase plate.}
\begin{enumerate}[label=\Alph*.]
\item Collide and break
\item Disappear
\item Overlap and interact
\item Fuse into a solid ring
\item Stop growing
\end{enumerate}
\textbf{Answer:} C \\
\textbf{Explanation:} This overlap is essential for cell elongation later.

\item \textbf{ATP provides \underline{\hspace{1cm}} for motor proteins to walk.}
\begin{enumerate}[label=\Alph*.]
\item Genetic information
\item Energy
\item Subunits
\item Oxygen
\item Heat
\end{enumerate}
\textbf{Answer:} B \\
\textbf{Explanation:} Walking is an active process requiring chemical energy.

\item \textbf{Cell plate formation starts with vesicles from the:}
\begin{enumerate}[label=\Alph*.]
\item Nucleus
\item Golgi apparatus
\item Mitochondria
\item Lysosome
\item Plasma membrane
\end{enumerate}
\textbf{Answer:} B \\
\textbf{Explanation:} Golgi vesicles carry cell wall materials to the center.

\item \textbf{Binary fission literally means:}
\begin{enumerate}[label=\Alph*.]
\item Mitosis
\item Division in half
\item Growth in size
\item DNA replication
\item Nuclear collapse
\end{enumerate}
\textbf{Answer:} B \\
\textbf{Explanation:} It refers to one cell splitting into two equal halves.

\item \textbf{Which protein related to tubulin helps pinch the bacterial plasma membrane?}
\begin{enumerate}[label=\Alph*.]
\item Actin-like protein
\item Cohesin
\item Tubulin-like protein
\item DNA polymerase
\item Myosin
\end{enumerate}
\textbf{Answer:} C \\
\textbf{Explanation:} Bacteria use proteins ancestral to tubulin for cytokinesis.

\item \textbf{In which phase of the cycle are liver cells typically found?}
\begin{enumerate}[label=\Alph*.]
\item Continuous M phase
\item G0 phase (reserve)
\item S phase
\item Anaphase
\item Prometaphase
\end{enumerate}
\textbf{Answer:} B \\
\textbf{Explanation:} They only divide to repair wounds or replace damage.

\item \textbf{Fusing an S phase cell with a G1 phase cell causes:}
\begin{enumerate}[label=\Alph*.]
\item Both cells to die
\item The G1 nucleus to immediately enter S phase
\item The S nucleus to revert to G1
\item No change
\item Cytokinesis
\end{enumerate}
\textbf{Answer:} B \\
\textbf{Explanation:} Molecules in the S phase cytoplasm signal the G1 nucleus to start DNA synthesis.

\item \textbf{Fusing an M phase cell with any interphase cell causes:}
\begin{enumerate}[label=\Alph*.]
\item The M cell to stop dividing
\item The interphase nucleus to enter mitosis
\item The formation of a single giant nucleus
\item Chromosomes to disappear
\item DNA replication
\end{enumerate}
\textbf{Answer:} B \\
\textbf{Explanation:} M phase cytoplasmic signals are very potent regulators.

\item \textbf{MPF activity peaks during which phase?}
\begin{enumerate}[label=\Alph*.]
\item G1
\item S
\item G2
\item M
\item G0
\end{enumerate}
\textbf{Answer:} D \\
\textbf{Explanation:} MPF triggers M phase, so its activity is highest then.

\item \textbf{Cyclin levels fall abruptly during:}
\begin{enumerate}[label=\Alph*.]
\item S phase
\item M phase
\item G1 phase
\item G2 phase
\item G0 phase
\end{enumerate}
\textbf{Answer:} B \\
\textbf{Explanation:} Degradation of cyclin inactivates MPF and ends the M phase.

\item \textbf{What is the "built-in clock" of the cell cycle?}
\begin{enumerate}[label=\Alph*.]
\item A stopwatch in the nucleus
\item Fluctuations in cell cycle control molecules
\item The sunrise and sunset
\item The speed of DNA replication
\item The size of the cell wall
\end{enumerate}
\textbf{Answer:} B \\
\textbf{Explanation:} Rhythmic levels of cyclins and Cdks pace the cycle.

\item \textbf{Cancer cells manage to \underline{\hspace{1cm}} the usual controls.}
\begin{enumerate}[label=\Alph*.]
\item Obey
\item Escape
\item Speed up
\item Strengthen
\item Ignore
\end{enumerate}
\textbf{Answer:} B \\
\textbf{Explanation:} Cancer is defined by unregulated, uncontrolled cell division.

\item \textbf{Which cells in a mature human do NOT divide at all?}
\begin{enumerate}[label=\Alph*.]
\item Skin cells
\item Liver cells
\item Fully formed nerve and muscle cells
\item Bone marrow cells
\item Intestinal lining cells
\end{enumerate}
\textbf{Answer:} C \\
\textbf{Explanation:} These are highly specialized cells that remain in G0 permanently.

\item \textbf{A multicellular organism starts as a:}
\begin{enumerate}[label=\Alph*.]
\item Two-celled embryo
\item Single cell (fertilized egg)
\item Cluster of genes
\item Spindle pole
\item Prokaryote
\end{enumerate}
\textbf{Answer:} B \\
\textbf{Explanation:} Every complex organism begins as a single zygote.

\item \textbf{Accuracy in DNA distribution is \underline{\hspace{1cm}} for life.}
\begin{enumerate}[label=\Alph*.]
\item Irrelevant
\item Crucial
\item Rare
\item Slow
\item Optional
\end{enumerate}
\textbf{Answer:} B \\
\textbf{Explanation:} Errors in DNA distribution lead to genetic instability and cell death.

\item \textbf{Which stage is conventionally excluded when describing the mitotic stages?}
\begin{enumerate}[label=\Alph*.]
\item Prophase
\item Anaphase
\item Cytokinesis (often considered separate/overlapping)
\item Prometaphase
\item Metaphase
\end{enumerate}
\textbf{Answer:} C \\
\textbf{Explanation:} Mitosis is the nuclear division; cytokinesis is the cytoplasmic division.

\item \textbf{Sister chromatids are often attached all along their lengths by \underline{\hspace{1cm}}.}
\begin{enumerate}[label=\Alph*.]
\item Glue
\item Cohesins
\item Spindle fibers
\item DNA polymerase
\item Asters
\end{enumerate}
\textbf{Answer:} B \\
\textbf{Explanation:} Cohesins ensure the identical copies stay together until anaphase.

\item \textbf{What is the rate of chromosome movement in micrometers per minute?}
\begin{enumerate}[label=\Alph*.]
\item 0.1
\item 1
\item 10
\item 100
\item 1000
\end{enumerate}
\textbf{Answer:} B \\
\textbf{Explanation:} They move at roughly 1 $\mu$m/min.

\item \textbf{What is the total length of DNA in a typical human cell compared to its diameter?}
\begin{enumerate}[label=\Alph*.]
\item 25 times
\item 250 times
\item 250,000 times
\item Equal
\item 1,000 times
\end{enumerate}
\textbf{Answer:} C \\
\textbf{Explanation:} 2 meters vs. micrometer scale diameters leads to a 250,000x difference.

\item \textbf{How many chromatids are in a cell with 6 duplicated chromosomes?}
\begin{enumerate}[label=\Alph*.]
\item 6
\item 12
\item 18
\item 24
\item 3
\end{enumerate}
\textbf{Answer:} B \\
\textbf{Explanation:} Each duplicated chromosome has two sister chromatids.

\item \textbf{The G2 checkpoint verifies that \underline{\hspace{1cm}} is complete.}
\begin{enumerate}[label=\Alph*.]
\item Cytokinesis
\item DNA replication
\item Prophase
\item Anaphase
\item G1
\end{enumerate}
\textbf{Answer:} B \\
\textbf{Explanation:} The cell must not start mitosis until its genome is fully copied.

\item \textbf{In binary fission, origins move toward \underline{\hspace{1cm}} ends.}
\begin{enumerate}[label=\Alph*.]
\item The same
\item Opposite
\item Circular
\item Interior
\item Random
\end{enumerate}
\textbf{Answer:} B \\
\textbf{Explanation:} This ensures each new cell gets one copy of the origin.

\item \textbf{Chromosomes are visible with a light microscope \underline{\hspace{1cm}}.}
\begin{enumerate}[label=\Alph*.]
\item Only during interphase
\item Only after they condense for division
\item Always
\item Never
\item Only in bacteria
\end{enumerate}
\textbf{Answer:} B \\
\textbf{Explanation:} They are too thin and diffuse as chromatin to see otherwise.

\item \textbf{The name "Mitosis" comes from the Greek word for:}
\begin{enumerate}[label=\Alph*.]
\item Nucleus
\item Thread
\item Division
\item Color
\item Body
\end{enumerate}
\textbf{Answer:} B \\
\textbf{Explanation:} (Note: while "chroma" is color and "soma" is body, "mitosis" relates to the thread-like appearance of chromosomes).

\item \textbf{A chromatid becomes a chromosome the moment:}
\begin{enumerate}[label=\Alph*.]
\item It is replicated
\item It enters the nucleus
\item Its sister chromatid is cleaved away in anaphase
\item It condenses in prophase
\item It enters the G0 phase
\end{enumerate}
\textbf{Answer:} C \\
\textbf{Explanation:} Once separate, they are independent chromosomal units.

\item \textbf{Which protein is a "motor" for walking?}
\begin{enumerate}[label=\Alph*.]
\item Tubulin
\item Actin
\item Dynein
\item Cohesin
\item Histone
\end{enumerate}
\textbf{Answer:} C \\
\textbf{Explanation:} Dynein and kinesin are the motor proteins of the spindle.

\item \textbf{The cell cycle control system is a \underline{\hspace{1cm}} operating set of molecules.}
\begin{enumerate}[label=\Alph*.]
\item Linearly
\item Cyclically
\item Randomly
\item Static
\item Slowly
\end{enumerate}
\textbf{Answer:} B \\
\textbf{Explanation:} It resets and repeats with every new cell generation.

\end{enumerate}
\end{document}